\chapter{How we reviewed the evidence}
\label{ch:evidence-method}

Suppose a reader finds a paper cited beside a product decision. The paper may
describe a failure, explain a mechanism, document a package, or suggest a
workflow worth funding. Those are four different jobs, and one search cannot
settle all four. This review keeps them separate so that the reader can see
what a source established, what we inferred, and what result would change the
recommendation.

The present document is not a completed systematic review. It is a living
evidence program with a frozen snapshot for this draft and a specified route to
independent completion. The reporting discipline follows PRISMA-S
\citep{rethlefsen2021prismas}; recent reviews of AI-assisted academic writing
and automated writing evaluation provide a topic-level coverage check
\citep{meng2026systematic,abudalfa2026awe}.

\section{Questions before sources}
\label{sec:method-questions}

Search terms are unstable when the question is not. The review therefore
starts with five product questions and one boundary question:

\begin{enumerate}
\item Do bounded writing or feedback interventions change author time,
      revision quality, or reader task completion?
\item Which properties of technical prose affect comprehension, trust, and
      the ability to reconstruct a decision?
\item How do authors use model suggestions, and what burden or error does the
      interaction introduce?
\item Which software components can represent, compare, and check a document
      while preserving equations, numbers, citations, and claims?
\item Which reader, privacy, disclosure, and integrity conditions constrain an
      acceptable workflow?
\item What evidence would distinguish an available component from an
      effective product bundle?
\end{enumerate}

The boundaries become visible in ordinary examples. A study of writing time
may inform the economic model but say nothing about technical comprehension. A
citation benchmark may resolve an identifier but say nothing about whether the
source supports a sentence. The claim ledger records both the question a source
answers and the questions it leaves open.

\begin{table}[H]
\centering
\small
\begin{tabularx}{\textwidth}{@{}p{0.22\textwidth}p{0.27\textwidth}L@{}}
\toprule
Question family & Preferred source types & Decision it can inform \\
\midrule
Observed failure & Internal case record with provenance; prospective reader
observation. & Whether a product problem is real in the target setting. \\
Intervention effect & Primary experiment, field study, or controlled
comparison with a named task and outcome. & Which outcome and comparator the
Humanvoice study should use. \\
Comprehension and genre & Technical-communication, readability, corpus, and
reader studies with population and task stated. & What a reader rubric and
genre baseline must include. \\
Interaction and judgment & Human--AI interaction studies and calibrated
evaluation research. & How to record suggestion, rejection, burden, and model
judge error. \\
Software capability & Official documentation, source, release, license, and
replayable fixture. & Whether a component belongs in the candidate inventory. \\
Integrity and policy & Standards, venue guidance, privacy and disclosure
research. & What the operating boundary and consent form must say. \\
\bottomrule
\end{tabularx}
\caption{The source hierarchy follows the decision question.}
\label{tab:method-question-sources}
\end{table}

\section{Search and retrieval protocol}
\label{sec:method-search}

The frozen snapshot follows four trails. The internal trail starts with repair
pairs, reader objections, build records, and the project's existing assets.
The scholarly trail searches generative assistance, writing feedback, technical
communication, readability, human--AI interaction, evaluation, disclosure,
provenance, and authorship detection. The software trail begins with the
capabilities the product needs and checks official repositories, documentation,
releases, and licenses. A fourth trail deliberately looks for trouble: null
effects, subgroup harms, detector false positives, citation failures, judge
bias, and packages that fail mathematical or privacy fixtures.

The protocol stores the query, source or database, date, time boundary, result
count, inclusion decision, and reason for exclusion. A query result is not an
included study. Known cornerstone items are marked as seeds and kept distinct
from items discovered by the search; otherwise a bibliography can appear to
have high recall simply because the desired papers were supplied in advance.

The first pass is broad enough to discover vocabulary. The second pass is
question-specific and records the exact title, abstract, or documentation
section inspected. Backward and forward citation chasing is reserved for
load-bearing sources. A source that changes only a peripheral example does not
receive the same effort as a source used to justify the primary endpoint or a
claim about harm.

Figure~\ref{fig:evidence-pipeline} shows the full path. The current audit has
completed the automated portions of the left-hand side; the independent and
full-text work in the middle remains a release condition for stronger claims.

\hvFigureEvidencePipeline

\begin{table}[H]
\centering
\small
\begin{tabularx}{\textwidth}{@{}p{0.20\textwidth}p{0.27\textwidth}L@{}}
\toprule
Pass & Activity & Output \\
\midrule
Vocabulary pass & Search broad terms and inspect reviews, reference lists,
and official capability pages. & Candidate terms, source identities, and
question assignments. \\
Question pass & Retrieve primary studies and documentation for each declared
question; record population, task, comparator, and outcome. & Claim records
with a preliminary source result and limit. \\
Challenge pass & Search for nulls, harms, subgroup variation, detector
failures, citation errors, and judge bias. & Counterevidence records and
revised claim classes. \\
Anchor pass & Inspect full text, methods, results, licenses, and release
information for load-bearing items. & Page or section anchors, appraisal, and
implementation constraints. \\
\bottomrule
\end{tabularx}
\caption{A staged search keeps discovery separate from support.}
\label{tab:method-passes}
\end{table}

The challenge pass is a protection against local optimization. An agent asked
only to find support for Humanvoice will preferentially retrieve papers that
make the product sound plausible. Requiring a counterevidence stream changes
the objective: a source is valuable when it narrows the claim as well as when
it supports one.

\section{Screening, appraisal, and adjudication}
\label{sec:method-screening}

Screening uses a small set of frozen reason codes. An item can be included as
direct evidence, included as adjacent evidence, retained as a challenge, or
excluded for wrong question, wrong population, insufficient source identity,
or inaccessible support. ``Interesting'' is not a reason code. The code must
explain why the item cannot currently change a product decision.

Two screeners should independently review the load-bearing sample. They record
the decision before seeing the other review; adjudication then resolves the
disagreement and preserves both original judgments. For software, the
equivalent is a second operator replaying the fixture from a clean environment.
One person can perform the first feasibility pass, but the result must not be
described as independent confirmation.

Appraisal is design-specific. For an intervention study, record assignment,
task, comparator, outcome construction, attrition, and heterogeneity. For a
reader study, record the reader role, time budget, rubric, and whether the
reader saw the source history. For a detector study, record the corpus,
editing, language, and false-positive analysis. For a package, record version,
license, input formats, network behavior, maintenance, and failure mode.

\begin{cardexample}{The same paper can receive two appraisals}
A field study may provide strong evidence that an assistant shortens email
completion time. It can be direct evidence for the question ``can assistance
save time on a bounded writing task?'' It is adjacent evidence for the question
``will a former professor accept a mathematical product proposal more quickly?''
The result enters both rows, but its fit and permitted inference differ. A
single undifferentiated quality label would hide that distinction.
\end{cardexample}

The current snapshot has completed the declared live retrieval checks, DOI
identity checks, and a bounded RePEc/IDEAS specialist import, but it has not
completed every independent screen or full-text anchor. The correct status is
therefore ``provisional'' for the survey and ``bounded'' for the product
recommendation. The missing work is listed as an implementation deliverable,
with an owner and a gate, rather than being silently converted into
confidence.

\section{From source result to product claim}
\label{sec:method-synthesis}

The synthesis record has five fields:

\begin{enumerate}
\item the source result, stated as close as possible to the inspected text;
\item the population, task, comparator, and outcome to which it belongs;
\item the limit or plausible alternative explanation;
\item the local implication for a Humanvoice object, fixture, or protocol; and
\item the claim class and next observation.
\end{enumerate}

This order matters. Starting with the product implication encourages a writer
to reinterpret a source until it sounds like a requirement. Starting with the
result and its limit makes the bridge inspectable.

\begin{table}[H]
\centering
\small
\begin{tabularx}{\textwidth}{@{}p{0.21\textwidth}L p{0.25\textwidth}@{}}
\toprule
Record field & Example of a disciplined entry & What it does not establish \\
\midrule
Source result & A bounded writing experiment reports lower completion time on
its stated task. & A benefit for expert technical review. \\
Limit & The task, population, and outcome differ from a proposal reader's
decision. & That the result is irrelevant. \\
Local implication & Measure author and reader time and declare the incumbent
bundle in the Humanvoice comparison. & That the measured effect will be
positive. \\
Claim class & Adjacent source result plus project hypothesis. & A validated
product claim. \\
Next observation & A held-out technical document with a timed reader rubric. &
That one feasibility case is enough. \\
\bottomrule
\end{tabularx}
\caption{A claim record preserves the bridge and its boundary.}
\label{tab:method-claim-record}
\end{table}

The matrix in Appendix~\ref{app:literature-matrix} is the expanded form of
this record. The main route uses prose because a decision-maker needs the
mechanism and implication together; the appendix makes it possible to audit
the rows without interrupting the argument.

\section{Software survey as an empirical review}
\label{sec:method-software}

Software changes faster than the scholarly literature, so a package survey
needs a version boundary. The current inventory records the version or commit
examined on the snapshot date, the official source, and whether the package
was actually run. A README-only row is an availability observation, not a
benchmark result. A smoke run is an installation observation, not an
operational pass. Only a held-out fixture and a human-burden measure can move
the candidate toward adoption.

The package review uses capability families rather than repository popularity:
source representation, build and cross-reference, protected comparison,
diagnostics, citation identity, local inference, and evaluation. For each
family, at least one positive and one deliberately failing fixture are
required. The failure fixture is important because a graceful abstention is a
product behavior, whereas a silent success is a safety defect.

A candidate's maintenance and license are part of technical fit. A package
with excellent parsing but an incompatible license may be useful as a
benchmark reference and unusable as a dependency. A package with a permissive
license but no recent release may be acceptable for a one-off migration and
unacceptable for a service. The disposition is therefore adopt, adapt, pilot,
or skip, with a reason that can be revisited when the version changes.

\section{Limits of the present review}
\label{sec:method-limits}

The method supports a transparent, challengeable product recommendation. It
does not yet support a claim that the literature has been exhaustively
searched, that every cited source has been independently verified, or that the
workflow improves expert writing. Those stronger claims require the gates
listed in the evidence status and the comparative program.

The method also cannot remove judgment. A reading brief, a claim scope, and
the usefulness of a finding remain partly adjudicated objects. The purpose of
the protocol is to expose where judgment enters, who owns it, and what
counterevidence could overturn it. That is a stronger scholarly posture than
pretending that a score has made the judgment disappear.

We can now ask what the sources actually say, then ask which software deserves
a fixture. That order keeps the bridge visible: first define how a source earns
a role, then interpret the result, then make a build choice. A reader can
disagree with any one choice without losing the path by which it was made.
